\section{现金目标：当前不能证明 3\% 净现金}

本案没有把任何历史产品分派当作未来现金预测。三个 ETF 的分配政策和历史记录不同；七只 REIT 都有“在符合条件时年度分配不低于可供分配金额 90\%”的政策，但这不是固定收益率。特别是，季度收入、费用、项目收现、扩募、保理和实际分派并不均匀，发行人已提示季度可分配金额不能简单年化。

\begin{exhibitbox}[表 13：历史现金观察不能替代 3\% 目标验证]
\scriptsize
\begin{tabularx}{\exhibitboxwidth}{L{1.25cm}L{2.75cm}L{3.25cm}X}
\toprule
\textbf{代码} & \textbf{可见历史现金资料} & \textbf{为何不能外推} & \textbf{进入净现金模型前仍需补齐} \\
\midrule
511010 & 2025 年现金分配 CNY21.042/10 份 & 单年度历史分配，不是票息承诺 & 实际结算、税费、申赎与未来分派资料 \\
510880 & 2023--2025 均有每 10 份分配记录 & 指数成份、利润、基金决定与税费会变化 & TTM 已结算现金、完整持仓、税后现金桥 \\
512890 & 年报称近三年未实施利润分配 & “红利低波”不代表基金持有人现金 & 基金分配政策、TTM 已结算现金与持仓质量 \\
508069--508096 & 年报历史分派率与 Q2 管理人年化口径 & Q2 未审计；季度现金、扩募、费用、保理和项目运营非线性 & 每项的可分配金额桥、实际结算、税费、债务/资本开支和压力情景 \\
\bottomrule
\end{tabularx}
\sourcenote{ETF 与 REIT 年报/Q2 核验包：\sourceid{S04--S25}。}
\end{exhibitbox}

合格的现金口径必须为滚动 12 个月的税后、费后、已结算、可复投现金；出售资产、未实现市值变动、NAV 增长、外部资金、借款以及未验证的返还本金均排除。若短缺，策略必须记录短缺，而不是以杠杆、低信用替代、被迫购入 REIT 或卖出本金修补。当前没有足够数据将任何产品的历史分派加总为组合 3\%。

\section{5\% 总回报目标：没有可复现实证就不是策略结果}

{\sloppy
总回报应以逐点产品集合和外部现金流分段的 TWRR 计算，现金只计入一次；投资者体验另以 MWRR 与 XIRR 报告。当前缺少历史成份/项目状态、分派、公司行为、再平衡、交易成本、税费、流动性/容量和基准序列，因而没有回放，也没有“每年 5\%”的证据。正式测试至少要与保护性流动性、期限匹配国债暴露和相关多资产替代方案比较，并报告回撤、命中率、换手、容量和现金短缺频率。\par}

\section{共同风险矩阵：名称不同，不等于独立}

\begin{exhibitbox}[表 14：跨产品共同风险与当前控制]
\small
\begin{tabularx}{\exhibitboxwidth}{L{2.3cm}L{4.0cm}X}
\toprule
\textbf{风险簇} & \textbf{成员与传导机制} & \textbf{当前控制/缺口} \\
\midrule
利率与估值 & 511010 的国债久期；所有 REIT 的融资成本和估值折现率 & 缺久期/DV01、同日 NAV/估值桥、债务到期表和联动压力测试。\\
红利/金融/周期权益 & 510880 的采矿/金融权重与 512890 的金融权重 & 需完整同日持仓、发行人集团重叠、因子回撤与现金质量分析。\\
租赁住房 & 508068、508077、508055：出租率、租金、续约、收缴、区域供给、政策/长期应付款 & 三只不得当作三份独立现金来源；需共同风险上限和项目级压力测试。\\
单资产或集中租户 & 508069 单公路、508080 最大租户 44.04\%、508011 前五租户 53.02\% & 需运营 KPI、租约/信用、续约、资本开支及债务资料。\\
能源/资源现金流 & 508096 的水光资源、电价、调度、补贴应收与保理 & 需发电、电价、应收、保理、借款和 NAV 口径勾稽。\\
市场可执行性 & 全部十项的点差、深度、容量、做市、账户和税务 & 单日成交额不等于可执行容量；所有 q 维持为零。\\
\bottomrule
\end{tabularx}
\sourcenote{ETF/REIT 穿透图谱和上交所快照：\sourceid{S04--S28}。}
\end{exhibitbox}

\section{由证据备忘录到实施的门禁表}

\begin{dashboardbox}[实施资格：全部阻断]
\begin{tabularx}{\textwidth}{L{3.3cm}X L{2.7cm}}
\textbf{回报证据} & 未完成固定产品集合的逐点 TWRR/已结算净现金回放与基准比较。 & \blocked \\
\textbf{风险预算} & 未批准共同因子、压力损失、集中度、流动性容量及回撤行动规则。 & \blocked \\
\textbf{适配性与账户} & 未提供账户、税务、结算、托管、投资者适配和资金用途资料。 & \blocked \\
\textbf{执行} & 缺同日 NAV/IOPV、点差、盘口深度、20 日容量、冲击成本和市场做市资料。 & \blocked \\
\textbf{产品数量 q} & 所有 ETF、REIT 与未命名类别均为零；不存在实时执行授权。 & \qzero \\
\end{tabularx}
\end{dashboardbox}

\section{下一次研究的优先级}

第一优先级不是追加候选，而是冻结两项未命名类别的有限产品清单，并完成十项产品的当前估值和执行资料。第二优先级是对 ETF 建立久期/当前持仓和现金质量包，对 REIT 建立资产级运营、债务、资本开支、同日估值和实际分派桥。第三优先级才是独立回放与风险预算审查。届时每一只产品仍需单独通过，不存在“策略整体通过即自动放行”的机制。
